% \documentclass[sigplan,screen,10pt]{acmart}
\documentclass[sigconf,review, anonymous]{acmart}
\settopmatter{printfolios=false,printccs=false,printacmref=false}
\renewcommand\footnotetextcopyrightpermission[1]{}

\usepackage{amsmath,amsfonts}
\usepackage[ruled,linesnumbered]{algorithm2e}
\usepackage{graphicx}
\usepackage{tabularx}
\usepackage{multirow}
\usepackage{multicol}
\usepackage{xcolor}
\usepackage{lipsum}
\usepackage{booktabs}
\usepackage{pifont}
\usepackage{xkeyval}
\usepackage{array}
\usepackage{subcaption}
\newcommand{\tabincell}[2]{\begin{tabular}{@{}#1@{}}#2\end{tabular}}
\usepackage{subcaption}
\captionsetup[sub]{labelfont=normalfont,textfont=normalfont}
\usepackage{microtype}
\usepackage{mathtools}

\bibliographystyle{ACM-Reference-Format}

%% Some recommended packages.
% \usepackage{booktabs}   %% For formal tables:
%                         %% http://ctan.org/pkg/booktabs
% \usepackage{subcaption} %% For complex figures with subfigures/subcaptions
%                         %% http://ctan.org/pkg/subcaption

%%% The following is specific to PPoPP '26 and the paper
%%% 'Accelerating Sparse Transformer Inference on GPU'
%%% by Wenhao Dai, Haodong Deng, Mengfei Rong, Xinyu Yang, Hongyu Liu, Fangxin Liu, Hailong Yang, Qianwen Cao, and Qingxiao Sun.
%%%
% \setcopyright{cc}
% \setcctype{by}
% \acmDOI{10.1145/3774934.3786434}
% \acmYear{2026}
% \copyrightyear{2026}
% \acmISBN{979-8-4007-2310-0/2026/01}
% \acmConference[PPoPP '26]{Proceedings of the 31st ACM SIGPLAN Annual Symposium on Principles and Practice of Parallel Programming}{January 31 -- February 4, 2026}{Sydney, NSW, Australia}
% \acmBooktitle{Proceedings of the 31st ACM SIGPLAN Annual Symposium on Principles and Practice of Parallel Programming (PPoPP '26), January 31 -- February 4, 2026, Sydney, NSW, Australia}
% \received{2025-09-01}
% \received[accepted]{2025-11-10}


\begin{document}

%% Title information
\title{MAGSAT: Memory-Access-Guided Fusion and Shape-Adaptive Auto-Tuning for Dynamic Tensor Programs}


%\author{Haodong Deng}
%\affiliation{
%  \department{Dept. of CST}
%  \institution{China University of Petroleum-Beijing} 
%  \city{Beijing}
%  \country{China}                    %% \country is recommended
%}
%\email{haodong.deng@student.cup.edu.cn}          %% \email is recommended
%
%\author{Jianuo Sheng}
%\affiliation{
%  \department{Dept. of CST}
%  \institution{China University of Petroleum-Beijing} 
%  \city{Beijing}
%  \country{China}                    %% \country is recommended
%}
%\email{jianuo.sheng@student.cup.edu.cn}          %% \email is recommended
%
%\author{Wenhao Dai}
%\affiliation{
%  \department{Dept. of CST}
%  \institution{China University of Petroleum-Beijing} 
%  \city{Beijing}
%  \country{China}                    %% \country is recommended
%}
%\email{wenhao.dai@student.cup.edu.cn}          %% \email is recommended
%
%
%\author{Hailong Yang}
%\affiliation{
%  \department{School of Computer Science and Engineering}   
%  \institution{Beihang University} 
%  \city{Beijing}
%  \country{China}
%}
%\email{hailong.yang@buaa.edu.cn}          %% \email is recommended
%
%
%\author{Qingxiao Sun}
%\authornote{Corresponding author} 
%\affiliation{
%  \department{School of Computer Science and Engineering}
%  \institution{Beihang University} 
%  \city{Beijing}
%  \country{China}
%}
%\email{qingxiaosun@buaa.edu.cn}

% \renewcommand{\shortauthors}{Deng et al.}


\begin{abstract}
% SQX：确定intro后再调整
\textcolor{red}{Large language models (LLMs) have driven increasing demand for high-performance GPU kernels. 
Many critical operators in these workloads exhibit dynamic shapes, while their performance is highly sensitive to shapes and kernel configurations. 
However, existing approaches typically rely on static fusion rules or shape-specific tuning results, which fail to generalize across diverse runtime shapes 
and often incur substantial performance loss under dynamic workloads. In addition, prior works rarely provide a unified mechanism to jointly decide 
fusion strategies and tuning configurations based on hardware-level performance characteristics.
To address these challenges, we propose MAGSAT, a unified framework for dynamic GPU operator optimization that combines memory-access-guided fusion and shape-adaptive auto-tuning. For fused operator generation, MAGSAT analyzes memory access behavior to determine whether and how operators should be fused, and maps different fusion decisions to corresponding compilation templates. For runtime optimization, MAGSAT employs a shape-adaptive auto-tuning strategy that integrates an analytical cost model with residual learning to efficiently select high-performance kernel configurations under unseen shapes. In this way, MAGSAT jointly optimizes fusion strategy, kernel template, and tuning parameters for dynamic workloads.
Experimental results show that MAGSAT consistently outperforms static fusion and static tuning baselines, achieving up to x speedup over existing methods.}
\end{abstract}


% \begin{CCSXML}
% <ccs2012>
%    <concept>
%        <concept_id>10010147.10010257</concept_id>
%        <concept_desc>Computing methodologies~Machine learning</concept_desc>
%        <concept_significance>500</concept_significance>
%        </concept>
%    <concept>
%        <concept_id>10010520.10010521.10010528.10010533</concept_id>
%        <concept_desc>Computer systems organization~Multiple instruction, single data</concept_desc>
%        <concept_significance>500</concept_significance>
%        </concept>
%  </ccs2012>
% \end{CCSXML}

% \ccsdesc[500]{Computing methodologies~Machine learning}
% \ccsdesc[500]{Computer systems organization~Multiple instruction, single data}


% \keywords{GPU, Sparse Transformer, Multi-head Attention, Operator Fusion}

\maketitle

%% --- Main Content ---
\section{Introduction}%1 page
\label{sec:introduction}

% 注意多加引用

% 第一段：深度学习应用的强大能力使其活跃于各个领域...从CV到NLP，再到大模型、多模态...范式在逐渐演化

% 第二段：深度学习的计算量很大，使其适合在GPU或领域特定加速器上部署...为了提高性能，厂商提出了cuDNN、xx等计算库..然而，开发高性能库非常耗时，且可能跟不上快速迭代的网络形态...

% 第三段：深度学习编译器支持张量程序的自动生成，且伴有调优器来最大化性能...然而其普遍针对静态形状...动态形状很常见...以Qwen模型为例...GEMM算子形状有几百个（具体数字）...针对每一种形状单独调优不现实...

% 第四段：动态形状编译器（给出学术界和工业界的典型例子）被提出...但其普遍只针对单独算子来调优...未考虑整网的算子融合机会...例如CI-MI之类的...只能取得次优性能...

% 第五段：近年来...基于tile的编译器逐渐流行如Triton, TileLang, cuTile等...这类的编译器在性能和灵活性上取得了权衡..方便用户编程...我们发现其在融合算子方面有很大潜力...调优后甚至取得比厂商算子库更高的性能...但同样其调优开销过于大...不适合于动态形状...

% 第六段：从以上的分析可知，优化端到端动态形状推理存在如下挑战：1）融合方案难以确定...2）缺乏高质量编译模板...3）调优成本太大...我们提出xxx...其首先...然后...最后...

% 第七段：据我们所知...这是第一份工作xxx....本文做出如下贡献....



\begin{itemize}

\item % 我们全面分析了动态形状下端到端全局推理的性能瓶颈并给出潜在的优化机会...

\item % 在图层，我们提出xx来确认融合方案并映射到编译模板...

\item % 在算子层，我们提出xx来针对指定形状快速确定参数配置...

\item % 我们实现了xx框架...其...实验结果表明....

\end{itemize} % 1.25 pages
\section{Background}%1 page
\label{sec:background}

\subsection{Dynamic Shape in Deep Learning} 
\label{sec:dynamic-shape}





\subsection{Tiled Program Compilation}
\label{sec:compilation}

%Triton, TileLang, cuTile ...


\subsection{Hierarchical Auto-tuning Techniques}
\label{sec:autotuning}

% 从算子融合到参数调优的分层自动调优技术，可借鉴STOF


   % 1.5 pages
\section{Motivation}%0.75 pages
\label{sec:motivation}

\subsection{Operator Fusion Opportunities}
\label{sec:fusion}

% 可做两部分实验，以GEMM为例，dynamic shape下，一个是单独算子，另一个是融合算子，单独算子时triton/tilelang能够比得过cublas/cudnn的比例较低；但融合算子对比cutlass等NVIDIA实现的融合算子，很大程度上会显示出优势？

% 从而服务于contribution 1-2，算子融合的必要性，以及triton/tilelang的模板保留了高性能和灵活性


\subsection{Variation in Optimal Configurations}
\label{sec:static}

% 以基于GEMM的融合算子为例，可给出热力图，证明不同形状下最优配置是有差别的，没有one-fits-all，服务于contribution 3

%\subsection{Tile Configuration Trade-offs}
%\label{sec:atuo-fusion}
%
%% 



\subsection{Bottleneck Shifts across Shapes}
\label{sec:bottleneck}

% 以GEMM和卷积算子为例，各画一个roofline model？不同shape下不同configuration的差别，证明比如随着shape变化，更趋向于memory-bound或compute-bound

   % 0.75 pages 
\section{Methodology} % 3 Pages  
\label{sec:method}


\subsection{Design Overview}
\label{sec:overview}

% 给出设计概要图，突出各个组件之间的关联关系


\subsection{Memory-Access-Driven Operator Fusion}
\label{sec:memory-access}

% 确定算子融合方案，其只和模型结构有关，在端到端推理过程中固定，和dynamic shape变化没有关系

\subsection{Template Mapping and Parallelization}
\label{sec:template}

% fused operator到template的映射，以及template的实现以及参数化供后续调优


\subsection{Parameter Configuration Determination}
\label{sec:determination}

% 用analytical model或者cost model？但引入cost model的话得有ablation study，以及对于训练cost model的数据集输入要详细介绍

\subsection{Implementation Details}
\label{sec:impdetails}
       % 3 pages
\section{Evaluation}  %3 pages
\label{sec:evaluation}

{\color{blue}

\subsection{Experiment Setup} 
\label{sec:setup}

\subsubsection{Hardware and Software Platforms} 

We evaluate MAGSAT on two GPU platforms: the NVIDIA A100 (Ampere architecture) and the H100 (Hopper architecture). 
These two platforms represent recent NVIDIA GPU generations with distinct differences in hardware resources, 
enabling us to assess whether the proposed framework remains effective across hardware generations. 
All experiments are conducted on machines running Ubuntu 22.04, with CUDA 12.6 and PyTorch 2.9.0.

\subsubsection{Comparison Configurations and Methods}

We evaluate both operator-level and model-level workloads. At the
operator level, we use representative GEMM and convolution workloads,
which cover the compute-intensive primitives commonly appearing in
Transformer and CNN models. At the model level, we evaluate four
representative neural networks: BERT~\cite{devlin2019bert},
LLaMA~\cite{touvron2023llama}, ResNet~\cite{he2016resnet}, and
AlexNet~\cite{krizhevsky2012alexnet}. These models exercise both
GEMM-dominated and Convolution-dominated computation patterns.

We compare MAGSAT with five baselines. \textbf{Torch} denotes eager-mode
PyTorch execution. \textbf{\texttt{torch.compile}} denotes the PyTorch 2.x
compiler stack based on TorchDynamo and TorchInductor~\cite{ansel2024pytorch}.
\textbf{CUTLASS} is NVIDIA's template-based GPU kernel library~\cite{nvidia22cutlass}.
\textbf{DietCode}~\cite{zheng2022dietcode} and
\textbf{Helix}~\cite{zhou2025helix} are state-of-the-art systems in dynamic-shape 
tensor program compilation. All experiments use FP16 inputs and outputs. For timing, we
perform 100 warm-up iterations before measurement to to record the average performance.



\subsection{Operator Performance}
\label{sec:operatorperf}

% 融合算子（单算子）的性能，如果单算子性能一般，就不用放

\subsection{End-to-end Performance}
\label{sec:end2end}

% 端到端网络推理的性能，LLM*2，CNN*2

\subsection{Ablation Study}
\label{sec:ablation}

% 组件的ablation测试，还需要再想想有啥能拆出来的；或者有什么框架特定的设死的参数，可做参数敏感性分析

\subsection{Parameter Tuning Efficiency}
\label{sec:tuning}

% 选到的配置相比全局最优配置的性能差别，以及在各种形状下取到最优配置的比例



\subsection{Overhead Analysis}
\label{sec:overhead}

% 主要是融合方案判定和参数调优的开销比例？

\subsection{Case Study: Qwen3-xxB Inference}
\label{sec:case}

% 可选（时间够用就加）：给个pretrained模型的端到端推理加速会更有说服力

}
   % 4.25 pages
% \vspace{-8pt}
\section{Related Work}  %0.25 pages
\label{sec:relatedwork}

{\color{blue}

% 写优化 
% 可分两段，从static shape算子优化，编译优化，全局（即端到端）优化；

% 过渡到dynamic shape算子优化，编译优化，全局优化；每段说明本文的优势在哪

\textbf{Static shapes.}
Static-shape deep learning optimization has been studied at the operator, compiler, and graph levels. Vendor libraries and template-based systems such as cuBLAS~\cite{nvidia24cublas}, cuDNN~\cite{chetlur2014cudnn}, CUTLASS~\cite{nvidia22cutlass}, Triton~\cite{tillet2019triton}, and TileLang~\cite{cheng2025pipethreader} provide highly optimized kernels for common tensor operators. Tensor compilers and auto-schedulers, including TVM~\cite{chen2018tvm}, AutoTVM~\cite{chen2018autotvm}, Ansor~\cite{zheng2020ansor}, and Halide~\cite{ragan2013halide}, automate schedule generation through loop transformations, cost models, and hardware-aware search. At the graph or end-to-end level, PyTorch~\cite{ansel2024pytorch}, TASO~\cite{jia2019taso}, Rammer~\cite{ma2020rammer}, DNNFusion~\cite{niu2021dnnfusion}, Welder~\cite{shi2023welder}, Chimera~\cite{zheng2023chimera}, Souffle~\cite{xia2024souffle}, and Mirage~\cite{wu2025mirage} exploit graph rewriting, operator fusion, global analysis, or multi-level superoptimization. These approaches typically assume fixed tensor shapes, static fusion rules, or shape-specific tuning results, making them less effective when runtime shapes change frequently.

\textbf{Dynamic shapes.}
Recent work has explored dynamic-shape workloads from both operator and compiler perspectives. At the operator level, DietCode~\cite{zheng2022dietcode}, MikPoly~\cite{yu2024mikpoly}, FTuner~\cite{mu2024ftuner}, and Helix~\cite{zhou2025helix} reduce dynamic-shape tuning or compilation overhead through shape-generic search spaces, runtime micro-kernel composition, uKernel-based tuning, or sample-free compilation. At the compiler level, Nimble~\cite{shen2021nimble}, DISC~\cite{zhu2021disc}, Relax~\cite{lai2025relax}, and CoRA~\cite{fegade2022cora} improve dynamic-shape execution through dispatching, symbolic shape representation, cross-level abstraction, or padding reduction. Unlike these works, our approach jointly reasons about memory-access-guided fusion, compilation template mapping, and shape-adaptive parameter tuning, enabling dynamic tensor programs to adapt both fusion decisions and kernel configurations according to runtime shape characteristics.

}
  % 0.5 pages
% \vspace{-8pt}
\section{Conclusion}%0.25 pages
\label{sec:conclusion}
   % 0.25 pages

% \section*{Acknowledgements}
% We sincerely thank our shepherd, Gagan Agrawal, and the anonymous reviewers for their insightful feedback that greatly improved this paper. This work is supported by National Natural Science Foundation of China (Grant No. 62402525, 62322201, U23B2020, 62402526), Beijing Natural Science Foundation (Grant No. 4244086), and CCF-Baidu Open Fund. Qingxiao Sun is the corresponding author.

 %and the Fundamental Research Funds for the Central Universities (Grant No. 2462023YJRC023)


% \balance
\bibliography{bibliography}



\end{document}
